% GENERATED FILE. Edit canonical lessons, then run npm run generate:books.
% canonical-source-hash: fnv1a64:7b78e819100a5062
% canonical-lessons: KA-C161-antisu, KA-C161-tilisu, KA-C161-bheti-madu, KA-C161-kaydirisu, KA-C161-raddu-madu

\chapter{To stick, To inform, To meet, To reserve, To cancel}
\label{ch:ka-to-stick-to-inform-to-meet-to-reserve-to-cancel}
\hlchaptermodality{161}

\begin{chapteropening}
\textbf{By the end of this chapter:} \emph{I can say to stick, to inform, to meet, to reserve and to cancel.}

Complete the last lesson of chapter 161: I can say to stick, to inform, to meet, to reserve and to cancel.
\end{chapteropening}

\section[\texorpdfstring{\kn{ಅಂಟಿಸು} (aṇṭisu)}{ಅಂಟಿಸು (aṇṭisu)}]{\kn{ಅಂಟಿಸು} (aṇṭisu) — to stick, to paste}

\label{lesson:KA-C161-antisu}



\begin{quote}
Before the new one: say the Kannada for to sneeze, then the Kannada for to shiver.
\end{quote}

\subsection*{You'll want to know: \kn{ಅಂಟಿಸು}}
\textbf{\kn{ಅಂಟಿಸು}} — \emph{aṇṭisu} — \textquotedblleft{}to stick, to paste\textquotedblright{}.

\begin{grammarlens}[title={What you've built}]
One more word for this chapter.
\end{grammarlens}

\subsection*{Your turn}
\begin{itemize}
\raggedright
  \item \emph{Say it:} \emph{aṇṭisu}
  \item \emph{Say it:} \emph{aṇṭisu}, once more
  \item \emph{Say it:} say \emph{aṇṭisu}
  \item \emph{Recall:} say the Kannada for to sneeze, then the Kannada for to shiver, then say \emph{aṇṭisu} again
\end{itemize}

\begin{tcolorbox}[breakable,colback=teal!4,colframe=teal!35!black,title={Before you move on}]
What does \kn{ಅಂಟಿಸು} mean? (\textquotedblleft{}To stick, to paste\textquotedblright{}.) Say it once more.
\end{tcolorbox}
\section[\texorpdfstring{\kn{ತಿಳಿಸು} (tiḷisu)}{ತಿಳಿಸು (tiḷisu)}]{\kn{ತಿಳಿಸು} (tiḷisu) — to inform, to let know}

\label{lesson:KA-C161-tilisu}



\begin{quote}
Before the new one: say the Kannada for to shiver, then the Kannada for to stick.
\end{quote}

\subsection*{You'll want to know: \kn{ತಿಳಿಸು}}
\textbf{\kn{ತಿಳಿಸು}} — \emph{tiḷisu} — \textquotedblleft{}to inform, to let know\textquotedblright{}.

\begin{grammarlens}[title={What you've built}]
One more word for this chapter.
\end{grammarlens}

\subsection*{Your turn}
\begin{itemize}
\raggedright
  \item \emph{Say it:} \emph{tiḷisu}
  \item \emph{Say it:} \emph{tiḷisu}, once more
  \item \emph{Say it:} say \emph{tiḷisu}
  \item \emph{Recall:} say the Kannada for to shiver, then the Kannada for to stick, then say \emph{tiḷisu} again
\end{itemize}

\begin{tcolorbox}[breakable,colback=teal!4,colframe=teal!35!black,title={Before you move on}]
What does \kn{ತಿಳಿಸು} mean? (\textquotedblleft{}To inform, to let know\textquotedblright{}.) Say it once more.
\end{tcolorbox}
\section[\texorpdfstring{\kn{ಭೇಟಿ} \kn{ಮಾಡು} (bhēṭi māḍu)}{ಭೇಟಿ ಮಾಡು (bhēṭi māḍu)}]{\kn{ಭೇಟಿ} \kn{ಮಾಡು} (bhēṭi māḍu) — to meet (someone)}

\label{lesson:KA-C161-bheti-madu}



\begin{quote}
Before the new one: say the Kannada for to stick, then the Kannada for to inform.
\end{quote}

\subsection*{You'll want to know: \kn{ಭೇಟಿ} \kn{ಮಾಡು}}
\textbf{\kn{ಭೇಟಿ} \kn{ಮಾಡು}} — \emph{bhēṭi māḍu} — \textquotedblleft{}to meet (someone)\textquotedblright{}.

\begin{grammarlens}[title={What you've built}]
One more word for this chapter.
\end{grammarlens}

\subsection*{Your turn}
\begin{itemize}
\raggedright
  \item \emph{Say it:} \emph{bhēṭi māḍu}
  \item \emph{Say it:} \emph{bhēṭi māḍu}, once more
  \item \emph{Say it:} say \emph{bhēṭi māḍu}
  \item \emph{Recall:} say the Kannada for to stick, then the Kannada for to inform, then say \emph{bhēṭi māḍu} again
\end{itemize}

\begin{tcolorbox}[breakable,colback=teal!4,colframe=teal!35!black,title={Before you move on}]
What does \kn{ಭೇಟಿ} \kn{ಮಾಡು} mean? (\textquotedblleft{}To meet (someone)\textquotedblright{}.) Say it once more.
\end{tcolorbox}
\section[\texorpdfstring{\kn{ಕಾಯ್ದಿರಿಸು} (kāydirisu)}{ಕಾಯ್ದಿರಿಸು (kāydirisu)}]{\kn{ಕಾಯ್ದಿರಿಸು} (kāydirisu) — to reserve, to book}

\label{lesson:KA-C161-kaydirisu}



\begin{quote}
Before the new one: say the Kannada for to inform, then the Kannada for to meet.
\end{quote}

\subsection*{You'll want to know: \kn{ಕಾಯ್ದಿರಿಸು}}
\textbf{\kn{ಕಾಯ್ದಿರಿಸು}} — \emph{kāydirisu} — \textquotedblleft{}to reserve, to book\textquotedblright{}.

\begin{grammarlens}[title={What you've built}]
One more word for this chapter.
\end{grammarlens}

\subsection*{Your turn}
\begin{itemize}
\raggedright
  \item \emph{Say it:} \emph{kāydirisu}
  \item \emph{Say it:} \emph{kāydirisu}, once more
  \item \emph{Say it:} say \emph{kāydirisu}
  \item \emph{Recall:} say the Kannada for to inform, then the Kannada for to meet, then say \emph{kāydirisu} again
\end{itemize}

\begin{tcolorbox}[breakable,colback=teal!4,colframe=teal!35!black,title={Before you move on}]
What does \kn{ಕಾಯ್ದಿರಿಸು} mean? (\textquotedblleft{}To reserve, to book\textquotedblright{}.) Say it once more.
\end{tcolorbox}
\section[\texorpdfstring{\kn{ರದ್ದು} \kn{ಮಾಡು} (raddu māḍu)}{ರದ್ದು ಮಾಡು (raddu māḍu)}]{\kn{ರದ್ದು} \kn{ಮಾಡು} (raddu māḍu) — to cancel}

\label{lesson:KA-C161-raddu-madu}



\begin{quote}
Before the new one: say the Kannada for to meet, then the Kannada for to reserve.
\end{quote}

\subsection*{You'll want to know: \kn{ರದ್ದು} \kn{ಮಾಡು}}
\textbf{\kn{ರದ್ದು} \kn{ಮಾಡು}} — \emph{raddu māḍu} — \textquotedblleft{}to cancel\textquotedblright{}.

\begin{grammarlens}[title={What you've built}]
One more word for this chapter.
\end{grammarlens}

\subsection*{Your turn}
\begin{itemize}
\raggedright
  \item \emph{Say it:} \emph{raddu māḍu}
  \item \emph{Say it:} \emph{raddu māḍu}, once more
  \item \emph{Say it:} say \emph{raddu māḍu}
  \item \emph{Recall:} say the Kannada for to meet, then the Kannada for to reserve, then say \emph{raddu māḍu} again
\end{itemize}

\begin{tcolorbox}[breakable,colback=teal!4,colframe=teal!35!black,title={Before you move on}]
What does \kn{ರದ್ದು} \kn{ಮಾಡು} mean? (\textquotedblleft{}To cancel\textquotedblright{}.) Say it once more.
\end{tcolorbox}
